\chapter[Lezione 4 --- 2 ottobre 2026]{Lezione 4: densità, intervalli ed estremi}
\label{chap:lezione-04}

\section{Densità dei numeri razionali}
\label{sec:lez04-densita-razionali}

La proprietà archimedea e la proprietà della parte intera, introdotte
rispettivamente nel \cref{thm:proprieta-archimedea,prop:parte-intera},
permettono di mostrare che ogni numero reale può essere approssimato
arbitrariamente da numeri razionali.

\begin{theorem}[Densità di \(\Q\) in \(\R\)]
\label{thm:densita-razionali}
Per ogni \(x\in\R\) e per ogni \(\varepsilon>0\) esiste \(z\in\Q\) tale che
\[
  z\le x<z+\varepsilon.
\]
\end{theorem}

\begin{proof}
Siano \(x\in\R\) ed \(\varepsilon>0\). Poiché
\[
  \frac{1}{\varepsilon}>0,
\]
per il \cref{thm:proprieta-archimedea} esiste \(q\in\N\) tale che
\[
  q>\frac{1}{\varepsilon}.
\]
In particolare \(q>0\).

Si considera allora il numero reale \(qx\). Per la
\cref{prop:parte-intera} esiste \(p\in\Z\) tale che
\[
  p\le qx<p+1.
\]
Dividendo per \(q>0\), il verso delle disuguaglianze non cambia e si
ottiene
\[
  \frac{p}{q}\le x<\frac{p+1}{q}
  =\frac{p}{q}+\frac{1}{q}.
\]

Dalla scelta di \(q\) segue
\[
  \frac{1}{q}<\varepsilon.
\]
Pertanto
\[
  \frac{p}{q}\le x
  <\frac{p}{q}+\frac{1}{q}
  <\frac{p}{q}+\varepsilon.
\]

Ponendo
\[
  z=\frac{p}{q},
\]
si ha \(z\in\Q\) e quindi
\[
  z\le x<z+\varepsilon.
\]
\end{proof}

\begin{figure}[htbp]
  \centering
  \begin{tikzpicture}[unipd/diagram]
  \fill[green!12] (0,-0.18) rectangle (4,0.18);
  \draw[->,unipd-muted] (-0.8,0) -- (5,0) node[right] {\(\R\)};
  \draw[very thick,unipd-primary] (0,0) -- (4,0);
  \fill[unipd-primary] (0,0) circle (2pt);
  \draw[unipd-primary,fill=white,thick] (4,0) circle (2pt);
  \fill[unipd-primary] (2.5,0) circle (2pt);
  \node[below=6pt] at (0,0) {\(z\)};
  \node[below=6pt] at (2.5,0) {\(x\)};
  \node[below=6pt] at (4,0) {\(z+\varepsilon\)};
  \draw[unipd-rule,thin] (0,0.2) -- (0,0.7) (4,0.2) -- (4,0.7);
  \draw[<->,unipd-muted] (0,0.6) -- (4,0.6)
    node[midway,above] {\(\varepsilon\)};
  \node at (2,-0.95) {\(z\le x<z+\varepsilon\)};
\end{tikzpicture}

  \caption{Interpretazione geometrica del \cref{thm:densita-razionali}: il
  numero razionale \(z\) dista da \(x\) meno di \(\varepsilon\).}
  \label{fig:densita-razionali}
\end{figure}

Un'ulteriore conseguenza del \cref{thm:densita-razionali}, insieme
all'esistenza di numeri irrazionali discussa nella
\cref{sec:lez03-completezza-radice-due}, è che tra due numeri reali
distinti esistono sia numeri razionali sia numeri irrazionali.

\begin{theorem}[Razionali e irrazionali tra due reali]
\label{thm:razionali-irrazionali-tra-reali}
Siano \(a,b\in\R\) con \(a<b\). Allora:
\begin{enumerate}
  \item esiste \(r\in\Q\) tale che
  \[
    a<r<b;
  \]
  \item esiste \(\alpha\in\R\setminus\Q\) tale che
  \[
    a<\alpha<b.
  \]
\end{enumerate}
\end{theorem}

La dimostrazione del \cref{thm:razionali-irrazionali-tra-reali} non viene
sviluppata in questa sede.

\section{Intervalli della retta reale}
\label{sec:lez04-intervalli}

Siano \(a,b\in\R\) con \(a<b\). Si introducono le seguenti notazioni:
\[
  (a,b)=\{x\in\R:a<x<b\},
\]
\[
  [a,b]=\{x\in\R:a\le x\le b\},
\]
\[
  [a,b)=\{x\in\R:a\le x<b\},
\]
\[
  (a,b]=\{x\in\R:a<x\le b\}.
\]

\begin{definition}[Intervallo della retta reale]
\label{def:intervallo-retta-reale}
Sia \(A\subseteq\R\), con \(A\neq\varnothing\). L'insieme \(A\) si dice
\emph{intervallo} se e solo se, per ogni \(x_1,x_2\in A\) con \(x_1<x_2\),
vale
\[
  \{x\in\R:x_1<x<x_2\}\subseteq A.
\]
\end{definition}

In altri termini, se due punti appartengono a un intervallo, allora anche
ogni punto reale compreso strettamente tra essi appartiene allo stesso
insieme.

\begin{example}
\label{ex:insieme-non-intervallo}
Si consideri
\[
  A=(0,1)\cup\{2\}.
\]
L'insieme \((0,1)\) è un intervallo, mentre \(A\) non lo è. Infatti si
possono scegliere, per esempio,
\[
  x_1=\frac12\in A,
  \qquad
  x_2=2\in A,
\]
ma
\[
  \frac32\in\left(\frac12,2\right)
  \qquad\text{e}\qquad
  \frac32\notin A.
\]
La condizione della \cref{def:intervallo-retta-reale} non è quindi
soddisfatta.
\end{example}

\begin{figure}[htbp]
  \centering
  \begin{tikzpicture}[unipd/diagram]
  \draw[->,unipd-muted] (-0.7,0) -- (5,0) node[right] {\(\R\)};
  \draw[very thick,unipd-primary] (0,0) -- (2,0);
  \foreach \p in {0,2}
    \draw[unipd-primary,fill=white,thick] (\p,0) circle (2pt);
  \foreach \p in {1,4}
    \fill[unipd-primary] (\p,0) circle (2pt);
  \draw[unipd-warning,fill=white,thick] (3,0) circle (2pt);
  \node[below=6pt] at (0,0) {\(0\)};
  \node[below=6pt] at (1,0) {\(\frac12\)};
  \node[below=6pt] at (2,0) {\(1\)};
  \node[below=6pt] at (3,0) {\(\frac32\)};
  \node[below=6pt] at (4,0) {\(2\)};
  \node[above=8pt] at (1,0) {\(\frac12\in A\)};
  \node[above=8pt,text=unipd-warning] at (3,0) {\(\frac32\notin A\)};
  \node[above=8pt] at (4.3,0) {\(2\in A\)};
  \node at (2,-1) {\(A=(0,1)\cup\{2\}\)};
\end{tikzpicture}

  \caption{L'insieme \(A=(0,1)\cup\{2\}\) non è un intervallo: tra
  \(\frac12\) e \(2\) esistono punti, come \(\frac32\), che non
  appartengono ad \(A\).}
  \label{fig:insieme-non-intervallo}
\end{figure}

\section{Maggioranti, minoranti e insiemi limitati}
\label{sec:lez04-maggioranti-minoranti}

Si consideri un insieme non vuoto \(A\subseteq\R\).

\begin{definition}[Maggiorante]
\label{def:maggiorante}
Un numero reale \(x\in\R\) si dice \emph{maggiorante} di \(A\) se e solo se
\[
  a\le x
  \qquad
  \forall a\in A.
\]
\end{definition}

Un maggiorante non deve necessariamente appartenere ad \(A\).

\begin{definition}[Minorante]
\label{def:minorante}
Un numero reale \(x\in\R\) si dice \emph{minorante} di \(A\) se e solo se
\[
  x\le a
  \qquad
  \forall a\in A.
\]
\end{definition}

Anche un minorante non deve necessariamente appartenere ad \(A\).

Per l'insieme considerato nell'\cref{ex:insieme-non-intervallo},
\[
  A=(0,1)\cup\{2\},
\]
il numero \(2\) è un maggiorante e anche \(3\) è un maggiorante. Analogamente,
\(0\) e \(-\frac12\) sono minoranti.

\begin{definition}[Limitatezza superiore e inferiore]
\label{def:limitatezza-superiore-inferiore}
Sia \(A\subseteq\R\), con \(A\neq\varnothing\).
\begin{enumerate}
  \item \(A\) si dice \emph{inferiormente limitato} se ammette almeno un
  minorante;
  \item \(A\) si dice \emph{superiormente limitato} se ammette almeno un
  maggiorante.
\end{enumerate}
\end{definition}

\begin{definition}[Insieme limitato]
\label{def:insieme-limitato}
Sia \(A\subseteq\R\), con \(A\neq\varnothing\). Si dice che \(A\) è
\emph{limitato} se è sia inferiormente sia superiormente limitato.
\end{definition}

L'insieme \(A=(0,1)\cup\{2\}\) è quindi limitato.

La limitatezza si conserva passando a un sottoinsieme non vuoto.

\begin{proposition}[Limitatezza dei sottoinsiemi]
\label{prop:limitatezza-sottoinsiemi}
Siano \(A,B\subseteq\R\), con \(B\neq\varnothing\) e \(B\subseteq A\).
\begin{enumerate}
  \item Se \(A\) è inferiormente limitato, allora \(B\) è inferiormente
  limitato.
  \item Se \(A\) è superiormente limitato, allora \(B\) è superiormente
  limitato.
\end{enumerate}
\end{proposition}

\begin{proof}
Si dimostra il primo punto. Se \(A\) è inferiormente limitato, esiste un
minorante \(x\in\R\) di \(A\), cioè
\[
  x\le a
  \qquad
  \forall a\in A.
\]
Poiché \(B\subseteq A\), per ogni \(b\in B\) vale anche
\[
  x\le b.
\]
Il numero \(x\) è dunque un minorante di \(B\), e quindi \(B\) è
inferiormente limitato.

Il secondo punto si dimostra in modo analogo sostituendo i minoranti con
i maggioranti.
\end{proof}

\subsection{Esempi fondamentali}
\label{subsec:lez04-esempi-limitatezza}

Per ogni \(n\in\N\) si ha
\[
  0\le n.
\]
Il numero \(0\) è quindi un minorante di \(\N\), e pertanto \(\N\) è
inferiormente limitato.

L'insieme \(\N\) non è invece superiormente limitato. Infatti, se esistesse
\(x\in\R\) tale che
\[
  n\le x
  \qquad
  \forall n\in\N,
\]
si avrebbe un maggiorante reale di \(\N\), in contraddizione con il
\cref{thm:proprieta-archimedea}.

Anche \(\Z\) non è superiormente limitato. Se lo fosse, poiché
\[
  \N\subseteq\Z,
\]
dalla \cref{prop:limitatezza-sottoinsiemi} seguirebbe che anche \(\N\)
sarebbe superiormente limitato, contro quanto appena dimostrato.

L'insieme \(\Z\) non è neppure inferiormente limitato. Supponendo per
assurdo che esista \(b\in\R\) tale che
\[
  b\le z
  \qquad
  \forall z\in\Z,
\]
moltiplicando per \(-1\) e applicando la
\cref{prop:compatibilita-ordine-interi} si ottiene
\[
  -b\ge -z
  \qquad
  \forall z\in\Z.
\]
Poiché, al variare di \(z\in\Z\), anche \(-z\) descrive tutti gli interi,
\(-b\) sarebbe un maggiorante di \(\Z\), in contraddizione con il fatto
che \(\Z\) non è superiormente limitato.

Poiché
\[
  \Z\subseteq\Q\subseteq\R,
\]
se \(\Q\) o \(\R\) fossero limitati superiormente oppure inferiormente,
la \cref{prop:limitatezza-sottoinsiemi} imporrebbe la corrispondente
limitatezza anche a \(\Z\). Ne segue che \(\Q\) e \(\R\) non sono né
superiormente né inferiormente limitati.

\section{Massimo e minimo}
\label{sec:lez04-massimo-minimo}

Nelle \cref{def:maggiorante,def:minorante} non si richiede che il
maggiorante o il minorante appartengano all'insieme. Per massimo e minimo,
invece, l'appartenenza è essenziale.

\begin{definition}[Massimo e minimo]
\label{def:massimo-minimo}
Sia \(A\subseteq\R\), con \(A\neq\varnothing\).
\begin{enumerate}
  \item \(A\) ammette \emph{massimo} se esiste \(x_0\in A\) tale che
  \[
    a\le x_0
    \qquad
    \forall a\in A.
  \]
  In tal caso si scrive \(x_0=\max A\).

  \item \(A\) ammette \emph{minimo} se esiste \(x_0\in A\) tale che
  \[
    x_0\le a
    \qquad
    \forall a\in A.
  \]
  In tal caso si scrive \(x_0=\min A\).
\end{enumerate}
\end{definition}

Si può quindi interpretare il massimo come un maggiorante che appartiene
all'insieme stesso e il minimo come un minorante che appartiene
all'insieme stesso.

Per esempio, se
\[
  A=(0,1],
\]
allora
\[
  \max A=1.
\]
L'insieme non ammette invece minimo. Infatti, per ogni \(a\in A\) si ha
\(0<a/2<a\le1\), dunque \(a/2\in A\) e \(a/2<a\); nessun elemento di
\(A\) può quindi essere il minimo.

Per descrivere in modo sistematico maggioranti e minoranti si introducono
i due insiemi
\[
  M_A:=\{x\in\R:x\text{ è un maggiorante di }A\},
\]
\[
  m_A:=\{x\in\R:x\text{ è un minorante di }A\}.
\]

Nel caso \(A=(0,1]\),
\[
  M_A=\{x\in\R:x\ge1\},
\]
mentre
\[
  m_A=\{x\in\R:x\le0\}.
\]
In particolare,
\[
  1=\min M_A
  \qquad\text{e}\qquad
  0=\max m_A.
\]

Si osserva quindi che \(M_A\) raccoglie tutti i maggioranti di \(A\),
mentre \(m_A\) raccoglie tutti i suoi minoranti. Nel caso considerato,
\(1\) è il più piccolo tra tutti i maggioranti e \(0\) è il più grande
tra tutti i minoranti.

\begin{theorem}[Estremi degli insiemi di maggioranti e minoranti]
\label{thm:estremi-maggioranti-minoranti}
Sia \(A\subseteq\R\), con \(A\neq\varnothing\).
\begin{enumerate}
  \item Se \(A\) è superiormente limitato, allora \(M_A\) ammette minimo.
  \item Se \(A\) è inferiormente limitato, allora \(m_A\) ammette massimo.
\end{enumerate}
\end{theorem}

La dimostrazione del \cref{thm:estremi-maggioranti-minoranti} non viene
sviluppata in questa sede.

\section{Estremo superiore ed estremo inferiore}
\label{sec:lez04-estremi}

Il \cref{thm:estremi-maggioranti-minoranti} permette di definire i due
estremi fondamentali di un insieme reale non vuoto.

\begin{definition}[Estremo inferiore]
\label{def:estremo-inferiore}
Sia \(A\subseteq\R\), con \(A\neq\varnothing\), inferiormente limitato.
Si chiama \emph{estremo inferiore} di \(A\), e si indica con
\(\inf A\), il massimo dell'insieme dei minoranti:
\[
  \inf A:=\max m_A.
\]
\end{definition}

\begin{definition}[Estremo superiore]
\label{def:estremo-superiore}
Sia \(A\subseteq\R\), con \(A\neq\varnothing\), superiormente limitato.
Si chiama \emph{estremo superiore} di \(A\), e si indica con
\(\sup A\), il minimo dell'insieme dei maggioranti:
\[
  \sup A:=\min M_A.
\]
\end{definition}

Per convenzione, se \(A\) non è superiormente limitato si scrive
\[
  \sup A=+\infty,
\]
mentre, se \(A\) non è inferiormente limitato, si scrive
\[
  \inf A=-\infty.
\]
I simboli \(+\infty\) e \(-\infty\) non rappresentano numeri reali;
sono utilizzati qui esclusivamente come convenzione notazionale.

Massimo e minimo, quando esistono, coincidono rispettivamente con estremo
superiore ed estremo inferiore.

\begin{proposition}[Relazione tra estremi, massimo e minimo]
\label{prop:relazione-estremi-massimo-minimo}
Sia \(A\subseteq\R\), con \(A\neq\varnothing\).
\begin{enumerate}
  \item Se \(\sup A\in A\), allora
  \[
    \sup A=\max A.
  \]
  \item Se \(\inf A\in A\), allora
  \[
    \inf A=\min A.
  \]
\end{enumerate}
\end{proposition}

La dimostrazione della \cref{prop:relazione-estremi-massimo-minimo} non viene
sviluppata in questa sede.

\section{Caratterizzazione di estremo superiore e inferiore}
\label{sec:lez04-caratterizzazione-estremi}

Le definizioni precedenti possono essere espresse mediante condizioni che
descrivono quanto gli elementi dell'insieme possano avvicinarsi ai suoi
estremi.

\begin{theorem}[Caratterizzazione dell'estremo superiore]
\label{thm:caratterizzazione-sup}
Sia \(A\subseteq\R\), con \(A\neq\varnothing\), superiormente limitato.
Per \(x\in\R\),
\[
  x=\sup A
\]
se e solo se valgono entrambe le condizioni
\[
  x\ge a
  \qquad
  \forall a\in A
\]
e
\[
  \forall\varepsilon>0\;\exists\bar a\in A
  \quad\text{tale che}\quad
  \bar a>x-\varepsilon.
\]
\end{theorem}

La prima condizione afferma, secondo la \cref{def:maggiorante}, che
\(x\) è un maggiorante di \(A\). La seconda garantisce che nessun numero
strettamente minore di \(x\) possa essere un maggiorante. Poiché ogni
elemento di \(A\) è minore o uguale a \(x\), le due condizioni insieme
equivalgono a richiedere che, per ogni \(\varepsilon>0\), esista
\(\bar a\in A\) tale che
\[
  x-\varepsilon<\bar a\le x.
\]

\begin{figure}[htbp]
  \centering
  \begin{tikzpicture}[unipd/diagram]
  \fill[green!12] (2,-0.18) rectangle (5,0.18);
  \draw[->,unipd-muted] (-0.5,0) -- (6,0) node[right] {\(\R\)};
  \foreach \p in {0,0.8,1.5,2.4,3.5,4.1,4.45}
    \fill[unipd-primary] (\p,0) circle (2pt);
  \node[above=8pt] at (0.8,0) {\(A\)};
  \node[fill=green!12,inner sep=1pt] at (4.7,0) {\(\cdots\)};
  \draw[unipd-muted] (2,-0.08) -- (2,0.08);
  \draw[unipd-muted] (5,-0.08) -- (5,0.08);
  \node[below=6pt] at (2,0) {\(x-\varepsilon\)};
  \node[below=6pt] at (5,0) {\(x=\sup A\)};
  \fill[green!40!black] (3.5,0) circle (2pt);
  \node[above=8pt,text=green!40!black] at (3.5,0) {\(\bar a\in A\)};
  \draw[unipd-rule,thin] (2,0.2) -- (2,0.85) (5,0.2) -- (5,0.85);
  \draw[<->,unipd-muted] (2,0.75) -- (5,0.75)
    node[midway,above] {\(\varepsilon>0\)};
  \node at (2.75,-1) {\(\forall\varepsilon>0\;\exists\bar a\in A:\quad x-\varepsilon<\bar a\le x\)};
\end{tikzpicture}

  \caption{Interpretazione della caratterizzazione di \(x=\sup A\): per
  ogni \(\varepsilon>0\) esiste \(\bar a\in A\) con
  \(x-\varepsilon<\bar a\le x\).}
  \label{fig:caratterizzazione-sup}
\end{figure}

\begin{proof}[Dimostrazione del \cref{thm:caratterizzazione-sup}]
Si dimostrano separatamente le due implicazioni.

\emph{Prima implicazione \((\Rightarrow)\).}
Si supponga
\[
  x=\sup A.
\]
Per la \cref{def:estremo-superiore},
\[
  x=\min M_A,
\]
quindi \(x\in M_A\). Per la \cref{def:maggiorante}, ciò equivale a
\[
  x\ge a
  \qquad
  \forall a\in A,
\]
e la prima condizione è verificata.

Sia ora \(\varepsilon>0\). Poiché
\[
  x-\varepsilon<x=\min M_A,
\]
si ha
\[
  x-\varepsilon\notin M_A.
\]
Per definizione di \(M_A\), un numero \(\alpha\in\R\) appartiene a \(M_A\)
se e solo se \(\alpha\ge a\) per ogni \(a\in A\). Pertanto
\(x-\varepsilon\notin M_A\) significa che esiste \(\bar a\in A\) tale che
\[
  x-\varepsilon<\bar a.
\]
Essendo \(\varepsilon>0\) arbitrario, anche la seconda condizione è
verificata.

\emph{Seconda implicazione \((\Leftarrow)\).}
Si supponga ora che valgano
\[
  x\ge a
  \qquad
  \forall a\in A
\]
e
\[
  \forall\varepsilon>0\;\exists\bar a\in A
  \quad\text{tale che}\quad
  x-\varepsilon<\bar a.
\]
La prima condizione implica, per la \cref{def:maggiorante}, che
\(x\in M_A\).

Poiché \(A\) è superiormente limitato, per il
\cref{thm:estremi-maggioranti-minoranti} l'insieme \(M_A\) ammette minimo.
Si pone
\[
  y:=\min M_A=\sup A.
\]
Dalla seconda condizione segue che, per ogni \(\varepsilon>0\),
\[
  x-\varepsilon\notin M_A.
\]
Infatti, per ciascun \(\varepsilon>0\) esiste \(\bar a\in A\) tale che
\(x-\varepsilon<\bar a\), mentre un maggiorante di \(A\) dovrebbe essere
maggiore o uguale a ogni elemento di \(A\).

Poiché \(y=\min M_A\), si ha \(y\in M_A\). Inoltre, ogni numero reale
\(t\ge y\) appartiene a \(M_A\): infatti, essendo \(y\) un maggiorante,
\[
  a\le y\le t
  \qquad
  \forall a\in A.
\]
Di conseguenza, da \(x-\varepsilon\notin M_A\) segue necessariamente
\[
  x-\varepsilon<y
  \qquad
  \forall\varepsilon>0.
\]

Se fosse \(y<x\), scegliendo
\[
  \varepsilon=x-y>0
\]
si otterrebbe
\[
  x-\varepsilon=y<y,
\]
che è assurdo. Deve quindi essere \(y\ge x\).

D'altra parte, essendo \(x\in M_A\) e \(y=\min M_A\), vale
\[
  y\le x.
\]
Pertanto \(x=y\). Poiché \(y=\sup A\), si conclude
\[
  x=\sup A.
\]
\end{proof}

\clearpage

\begin{theorem}[Caratterizzazione dell'estremo inferiore]
\label{thm:caratterizzazione-inf}
Sia \(A\subseteq\R\), con \(A\neq\varnothing\), inferiormente limitato.
Per \(x\in\R\),
\[
  x=\inf A
\]
se e solo se valgono entrambe le condizioni
\[
  x\le a
  \qquad
  \forall a\in A
\]
e
\[
  \forall\varepsilon>0\;\exists\bar a\in A
  \quad\text{tale che}\quad
  \bar a<x+\varepsilon.
\]
\end{theorem}

La prima condizione afferma, secondo la \cref{def:minorante}, che
\(x\) è un minorante di \(A\). La seconda garantisce che nessun numero
strettamente maggiore di \(x\) possa essere un minorante. Poiché ogni
elemento di \(A\) è maggiore o uguale a \(x\), le due condizioni insieme
equivalgono a richiedere che, per ogni \(\varepsilon>0\), esista
\(\bar a\in A\) tale che
\[
  x\le\bar a<x+\varepsilon.
\]

\begin{figure}[htbp]
  \centering
  \begin{tikzpicture}[unipd/diagram]
  \fill[green!12] (0,-0.18) rectangle (3,0.18);
  \draw[->,unipd-muted] (-1,0) -- (5.5,0) node[right] {\(\R\)};
  \foreach \p in {0.55,0.9,1.5,2.6,3.5,4.2,5}
    \fill[unipd-primary] (\p,0) circle (2pt);
  \node[above=8pt] at (4.2,0) {\(A\)};
  \node[fill=green!12,inner sep=1pt] at (0.3,0) {\(\cdots\)};
  \draw[unipd-muted] (0,-0.08) -- (0,0.08);
  \draw[unipd-muted] (3,-0.08) -- (3,0.08);
  \node[below=6pt] at (0,0) {\(x=\inf A\)};
  \node[below=6pt] at (3,0) {\(x+\varepsilon\)};
  \fill[green!40!black] (1.5,0) circle (2pt);
  \node[above=8pt,text=green!40!black] at (1.5,0) {\(\bar a\in A\)};
  \draw[unipd-rule,thin] (0,0.2) -- (0,0.85) (3,0.2) -- (3,0.85);
  \draw[<->,unipd-muted] (0,0.75) -- (3,0.75)
    node[midway,above] {\(\varepsilon>0\)};
  \node at (2.25,-1) {\(\forall\varepsilon>0\;\exists\bar a\in A:\quad x\le\bar a<x+\varepsilon\)};
\end{tikzpicture}

  \caption{Interpretazione della caratterizzazione di \(x=\inf A\): per
  ogni \(\varepsilon>0\) esiste \(\bar a\in A\) con
  \(x\le\bar a<x+\varepsilon\).}
  \label{fig:caratterizzazione-inf}
\end{figure}


La dimostrazione del \cref{thm:caratterizzazione-inf} viene lasciata al
lettore come esercizio, in modo da ripercorrere in forma simmetrica il
ragionamento appena svolto.

\begin{exercise}[Dimostrazione della caratterizzazione dell'estremo inferiore]
\label{ex:caratterizzazione-inf}
Dimostrare il \cref{thm:caratterizzazione-inf}. Si suggerisce di seguire i
due versi della dimostrazione del \cref{thm:caratterizzazione-sup},
sostituendo all'insieme \(M_A\) dei maggioranti l'insieme \(m_A\) dei
minoranti e utilizzando
\[
  \inf A=\max m_A.
\]
\end{exercise}

Dopo aver svolto autonomamente l'esercizio, si può confrontare il proprio
ragionamento con la seguente soluzione.

\begin{proof}[Soluzione dell'esercizio]
Si dimostrano separatamente le due implicazioni.

\emph{Prima implicazione \((\Rightarrow)\).}
Si supponga
\[
  x=\inf A.
\]
Per la \cref{def:estremo-inferiore},
\[
  x=\max m_A,
\]
quindi \(x\in m_A\). Per la \cref{def:minorante}, ciò equivale a
\[
  x\le a
  \qquad
  \forall a\in A,
\]
e la prima condizione è verificata.

Sia ora \(\varepsilon>0\). Poiché
\[
  x+\varepsilon>x=\max m_A,
\]
si ha
\[
  x+\varepsilon\notin m_A.
\]
Per definizione di \(m_A\), un numero \(\alpha\in\R\) appartiene a \(m_A\)
se e solo se \(\alpha\le a\) per ogni \(a\in A\). Pertanto
\(x+\varepsilon\notin m_A\) significa che esiste \(\bar a\in A\) tale che
\[
  \bar a<x+\varepsilon.
\]
Essendo \(\varepsilon>0\) arbitrario, anche la seconda condizione è
verificata.

\emph{Seconda implicazione \((\Leftarrow)\).}
Si supponga ora che valgano
\[
  x\le a
  \qquad
  \forall a\in A
\]
e
\[
  \forall\varepsilon>0\;\exists\bar a\in A
  \quad\text{tale che}\quad
  \bar a<x+\varepsilon.
\]
La prima condizione implica, per la \cref{def:minorante}, che
\(x\in m_A\).

Poiché \(A\) è inferiormente limitato, per il
\cref{thm:estremi-maggioranti-minoranti} l'insieme \(m_A\) ammette massimo.
Si pone
\[
  y:=\max m_A=\inf A.
\]
Dalla seconda condizione segue che, per ogni \(\varepsilon>0\),
\[
  x+\varepsilon\notin m_A.
\]
Infatti, per ciascun \(\varepsilon>0\) esiste \(\bar a\in A\) tale che
\(\bar a<x+\varepsilon\), mentre un minorante di \(A\) dovrebbe essere
minore o uguale a ogni elemento di \(A\).

Poiché \(y=\max m_A\), si ha \(y\in m_A\). Inoltre, ogni numero reale
\(t\le y\) appartiene a \(m_A\): infatti, essendo \(y\) un minorante,
\[
  t\le y\le a
  \qquad
  \forall a\in A.
\]
Di conseguenza, da \(x+\varepsilon\notin m_A\) segue necessariamente
\[
  y<x+\varepsilon
  \qquad
  \forall\varepsilon>0.
\]

Se fosse \(y>x\), scegliendo
\[
  \varepsilon=y-x>0
\]
si otterrebbe
\[
  y<x+\varepsilon=y,
\]
che è assurdo. Deve quindi essere \(y\le x\).

D'altra parte, essendo \(x\in m_A\) e \(y=\max m_A\), vale
\[
  x\le y.
\]
Pertanto \(x=y\). Poiché \(y=\inf A\), si conclude
\[
  x=\inf A.
\]
\end{proof}
